\chapter{Collaboration}

    \section{Industry Collaboration}

        \begin{figure}[h]
            \includegraphics[width=0.25\linewidth]{Photos/neo4j.png}
        \end{figure}

        \paragraph{} Since the early stages of our Honours Thesis planning, we engaged with \textbf{neo4j} and \textbf{Dr. Jim Webber} to explore opportunities for a strategic research collaboration. By Week 3, this interaction formalised into an official partnership established to support both the academic components of our thesis and our case study.
        
        \paragraph{} \textbf{neo4j} is a leading native graph database management system designed for modelling, querying, and analysing highly connected data. The property graph model, optimised traversal engine, and mature eco-system of tools make this well suited for representing complex technical domains, regulatory structures, and system inter-dependencies.\textbf{ neo4j’s} capabilities in graph analytics, constraint reasoning, and knowledge representation provide a robust foundation for developing explainable, queryable engineering knowledge bases, particularly where hierarchical standards, multi-layered requirements, and cross-disciplinary design artefacts must be integrated.
        
        \paragraph{} Through this partnership, we are leveraging \textbf{neo4j’s} technology and expert guidance to construct an ontology-driven knowledge graph that addresses Sunswift’s multifaceted engineering, compliance, and traceability challenges. By embedding structured regulatory knowledge into a graph-based systems engineering workflow, the collaboration enables us to investigate scalable, knowledge-driven approaches for aligning complex automotive subsystem design with the Australian Design Rules (ADR's).
        
        \paragraph{} This partnership positions our research as an exemplar of industry–academic collaboration. Our goal is to demonstrate how graph technologies can support regulatory-aligned systems engineering at scale, producing outputs that can be published, disseminated, and showcased as a model for future research and industry adoption.

    \section{Academic Collaboration}

        \paragraph{} We were referred to \textbf{Mohammed Malabari} by \textbf{Dr. Shiva Abdoli}, who directed us toward her research group for technical alignment and deeper expertise in knowledge representation. Mohammed, a PhD researcher within Dr. Abdoli’s group, specialises in machine learning and possesses significant experience in knowledge graph development and reasoning.
        
        \paragraph{} Our meeting with him served as an important technical checkpoint, enabling us to clarify several aspects of our implementation strategy. In particular, we examined whether a recommender system was required, explored the potential integration of a MCP, and assessed the overall feasibility and coherence of our proposed approach. This engagement ensured that each methodological pathway we intended to pursue contributed meaningful value to the project’s objectives and strengthened the technical rigour of our research outcomes.